\topic{Abbildung A.9: Warum bedeutet eins minus neu durch alt eine Einsparung?}
\question{Was bedeuten die Funktionsausdrücke in Gleichung 7.10?}
$N_{\mathrm{MAC}}(H)$ bezeichnet die Anzahl der MAC-Operationen bei der Hidden Size $H$. Die Klammer ist das Argument der Funktion, keine zusätzliche Multiplikation mit $N_{\mathrm{MAC}}$.
\[
\underbrace{N_{\mathrm{MAC}}(H)}_{\text{alter Aufwand}},\qquad
\underbrace{N_{\mathrm{MAC}}((1-p)H)}_{\text{neuer Aufwand}}.
\]
$p$ ist der entfernte Anteil der Hidden-Kanäle. Deshalb ist $(1-p)H$ die verbleibende Hidden Size. $D$ (Eingangsbreite) und $M$ (MLP-Breite) bleiben hier unverändert.

\question{Wie wird daraus alt minus neu, geteilt durch alt?}
Ausgangspunkt ist die Formel aus der Dissertation:
\[
R_{\mathrm{MAC}}=1-\frac{\text{neu}}{\text{alt}}.
\]
\textbf{Schritt 1:} Weil der alte Aufwand positiv ist, dürfen wir die Eins als gleichwertigen Bruch schreiben:
\[
1=\frac{\text{alt}}{\text{alt}}.
\]
\textbf{Schritt 2:} Diesen Bruch einsetzen:
\[
R_{\mathrm{MAC}}=\frac{\text{alt}}{\text{alt}}-\frac{\text{neu}}{\text{alt}}.
\]
\textbf{Schritt 3:} Bei gleichem Nenner werden die Zähler subtrahiert:
\[
\boxed{R_{\mathrm{MAC}}=\frac{\text{alt}-\text{neu}}{\text{alt}}.}
\]
Mit den Modellbezeichnungen:
\[
1-\frac{N_{\mathrm{MAC}}((1-p)H)}{N_{\mathrm{MAC}}(H)}
=\frac{N_{\mathrm{MAC}}(H)-N_{\mathrm{MAC}}((1-p)H)}{N_{\mathrm{MAC}}(H)}.
\]
Es wird nichts zusätzlich abgezogen: Beide Schreibweisen sind mathematisch identisch. Der Zähler rechts ist die eingesparte Anzahl an Operationen; die Division durch den alten Aufwand macht daraus einen relativen Anteil.

\question{Ein Zahlenbeispiel zum Merken}
Vorher 100 Operationen, nachher 60:
\[
1-\frac{60}{100}=\frac{100}{100}-\frac{60}{100}
=\frac{100-60}{100}=0{,}4=40\,\%.
\]
\textbf{Merksatz:} Neu durch alt ist der verbleibende Anteil. Eins minus diesen Anteil ist die Einsparung. Für den Prozentwert wird das Ergebnis mit 100 multipliziert.

\topic{Gleichung 7.10: Vollständige Herleitung}
\question{1. Ausgangspunkt: Welche Operationen werden gezählt?}
$H$ bezeichnet die ursprüngliche LSTM-Hidden-Size, $D$ die Anzahl der Eingangsfeatures, $M$ die MLP-Hidden-Breite und $p$ den entfernten Kanalanteil. $D$ und $M$ bleiben beim betrachteten Pruning konstant.
\[
N_{\mathrm{MAC}}(H)=\underbrace{4H(D+H)}_{\text{LSTM}}+
\underbrace{HM+M}_{\text{MLP}}.
\]
Das ist Gleichung 7.9: vier LSTM-Gates, die Abbildung vom Hidden State auf den MLP und dessen skalarer Ausgang. Bias-Additionen, Aktivierungsfunktionen und Speicherzugriffe werden hier nicht als MACs gezählt.
\[
N_{\mathrm{MAC}}(H)=4HD+4H^2+HM+M
=\boxed{4H^2+(4D+M)H+M}.
\]
\question{2. Neue Hidden Size einsetzen}
Nach Entfernen des Anteils $p$ bleibt $H'=(1-p)H$. Jedes $H$ in der MAC-Formel wird dadurch ersetzt:
\[
\begin{aligned}
N_{\mathrm{MAC}}(H')
&=4\bigl((1-p)H\bigr)^2+(4D+M)(1-p)H+M\\
&=4(1-p)^2H^2+(4D+M)(1-p)H+M.
\end{aligned}
\]
\question{3. Alten und neuen Aufwand in die Einsparung einsetzen}
\[
R(H,p)=1-\frac{N_{\mathrm{MAC}}(H')}{N_{\mathrm{MAC}}(H)}
=\frac{N_{\mathrm{MAC}}(H)-N_{\mathrm{MAC}}(H')}{N_{\mathrm{MAC}}(H)}.
\]
Ausgeschrieben:
\[
R(H,p)=\frac{4H^2+(4D+M)H+M-\left[4(1-p)^2H^2+(4D+M)(1-p)H+M\right]}
{4H^2+(4D+M)H+M}.
\]
\question{4. Den Zähler nach gleichen Termen gruppieren}
\[
\begin{aligned}
\text{Zähler}={}&4H^2\left[1-(1-p)^2\right]\\
&+(4D+M)H\left[1-(1-p)\right]+(M-M).
\end{aligned}
\]
Die drei Klammern werden einzeln vereinfacht:
\[
1-(1-p)^2=1-(1-2p+p^2)=2p-p^2,
\]
\[
1-(1-p)=p,\qquad M-M=0.
\]
\textbf{Ergebnis: Gleichung 7.10}
\[
\boxed{R(H,p)=\frac{4(2p-p^2)H^2+p(4D+M)H}{4H^2+(4D+M)H+M}.}
\]

\topic{Gleichung 7.10: Grenzwert und konkrete Modelle}
\question{5. Warum ergibt sich für große Hidden Sizes der Grenzwert $2p-p^2$?}
Wir teilen Zähler und Nenner durch $H^2$:
\[
R(H,p)=\frac{4(2p-p^2)+\dfrac{p(4D+M)}{H}}
{4+\dfrac{4D+M}{H}+\dfrac{M}{H^2}}.
\]
Für $H\to\infty$ bei festem $D$, $M$ und $p$ verschwinden die Terme mit $H$ oder $H^2$ im Nenner:
\[
\lim_{H\to\infty}R(H,p)=\frac{4(2p-p^2)}4=\boxed{2p-p^2}.
\]
Der MLP ist in der vollständigen Rechnung enthalten. Der Grenzwert entsteht, weil der quadratische rekurrente LSTM-Anteil schließlich dominiert. Der Eingangsteil und der MLP wachsen bei festen Breiten nur linear mit $H$.

\question{6. Warum sind es bei 30 Prozent Pruning theoretisch 51 Prozent?}
\[
p=0{,}30\quad\Rightarrow\quad 2p-p^2=0{,}60-0{,}09=0{,}51=51\,\%.
\]
Anschaulich bleiben bei einer rekurrenten Matrix 70 Prozent der Zeilen und 70 Prozent der Spalten: $0{,}7\cdot0{,}7=0{,}49$. Es bleiben 49 Prozent dieses quadratischen Anteils, also werden 51 Prozent eingespart. Beim Gesamtmodell ist das der Grenzwert, nicht die tatsächliche Einsparung jeder endlichen Architektur.

\question{7. Was ergibt sich für eure tatsächlich verwendeten Größen?}
Beide Modelle haben $D=6$. Beim SOC ist $M=64$, beim SOH $M=128$. Die ganzzahligen Hidden Sizes sind SOC $64\to45$ und SOH $128\to90$, also jeweils etwa 30 Prozent Pruning.
\[
\begin{aligned}
N_{\mathrm{SOC,alt}}&=4\cdot64(6+64)+64\cdot64+64=22\,080,\\
N_{\mathrm{SOC,neu}}&=4\cdot45(6+45)+45\cdot64+64=12\,124,\\
R_{\mathrm{SOC}}&=1-\frac{12\,124}{22\,080}\approx45{,}1\,\%.
\end{aligned}
\]
\[
\begin{aligned}
N_{\mathrm{SOH,alt}}&=4\cdot128(6+128)+128\cdot128+128=85\,120,\\
N_{\mathrm{SOH,neu}}&=4\cdot90(6+90)+90\cdot128+128=46\,208,\\
R_{\mathrm{SOH}}&=1-\frac{46\,208}{85\,120}\approx45{,}7\,\%.
\end{aligned}
\]
\textbf{Abgrenzung:} $4H(D+H+1)$ ist die LSTM-Parameterzahl bei einem Biasvektor pro Gate, nicht die hier verwendete MAC-Zahl des gesamten Modells. Der MLP fehlt darin. In Abbildung A.9 liegen die MAC-Kurven bei gleichem $H$ verschieden, weil die MLP-Breite beim SOH größer bleibt. Eine MAC-Einsparung ist außerdem nicht automatisch die gleiche prozentuale Laufzeiteinsparung.
